\documentclass[12pt,a4paper]{article}
\usepackage[UTF8]{ctex}
\usepackage{amsmath,amssymb,amsthm}
\usepackage{geometry}
\usepackage{bm}

\geometry{margin=2.5cm}

\title{向量平行的定义与性质：为什么标量倍数意味着平行？}
\author{第8章补充说明}
\date{}

\begin{document}

\maketitle

\section{问题提出}

\textbf{问题}：为什么如果角动量$\bm{L}$是角速度$\bm{\omega}$的标量倍数（即$\bm{L} = k \bm{\omega}$，其中$k$是标量），那么两者平行？

\section{向量平行的定义}

\subsection{数学定义}

\textbf{定义}：两个非零向量$\bm{a}$和$\bm{b}$平行，当且仅当存在标量$k \neq 0$使得：
\begin{equation}
\bm{a} = k \bm{b}
\label{eq:parallel_definition}
\end{equation}

或者等价地：
\begin{equation}
\bm{b} = \frac{1}{k} \bm{a} = k' \bm{a}
\end{equation}

其中$k' = \frac{1}{k}$。

\subsection{等价表述}

向量$\bm{a}$和$\bm{b}$平行的等价表述：
\begin{enumerate}
\item 存在标量$k$使得$\bm{a} = k \bm{b}$
\item 两个向量方向相同或相反（相差一个标量因子）
\item 两个向量的叉积为零：$\bm{a} \times \bm{b} = \bm{0}$
\item 两个向量的夹角为$0$或$\pi$（即$\cos\theta = \pm 1$）
\end{enumerate}

\section{为什么标量倍数意味着平行？}

\subsection{几何直观}

\textbf{几何意义}：
\begin{itemize}
\item 如果$\bm{a} = k \bm{b}$，那么$\bm{a}$和$\bm{b}$指向相同或相反的方向
\item 标量$k$只改变向量的长度（大小），不改变方向
\item 如果$k > 0$，两者方向相同
\item 如果$k < 0$，两者方向相反
\item 如果$k = 0$，则$\bm{a} = \bm{0}$（零向量与任何向量平行）
\end{itemize}

\subsection{数学证明}

\textbf{证明}：如果$\bm{a} = k \bm{b}$，那么$\bm{a}$和$\bm{b}$平行。

\textbf{方法1：使用定义}

根据向量平行的定义，如果存在标量$k$使得$\bm{a} = k \bm{b}$，那么$\bm{a}$和$\bm{b}$平行。这是定义本身，因此直接成立。

\textbf{方法2：使用叉积}

如果$\bm{a} = k \bm{b}$，那么：
\begin{align}
\bm{a} \times \bm{b} &= (k \bm{b}) \times \bm{b} \\
&= k (\bm{b} \times \bm{b}) \\
&= k \cdot \bm{0} \\
&= \bm{0}
\end{align}

由于$\bm{a} \times \bm{b} = \bm{0}$，因此$\bm{a}$和$\bm{b}$平行。

\textbf{方法3：使用点积和夹角}

两个向量的夹角$\theta$满足：
\begin{equation}
\cos\theta = \frac{\bm{a} \cdot \bm{b}}{\|\bm{a}\| \|\bm{b}\|}
\end{equation}

如果$\bm{a} = k \bm{b}$，那么：
\begin{align}
\bm{a} \cdot \bm{b} &= (k \bm{b}) \cdot \bm{b} \\
&= k (\bm{b} \cdot \bm{b}) \\
&= k \|\bm{b}\|^2
\end{align}

同时：
\begin{equation}
\|\bm{a}\| = \|k \bm{b}\| = |k| \|\bm{b}\|
\end{equation}

因此：
\begin{align}
\cos\theta &= \frac{k \|\bm{b}\|^2}{|k| \|\bm{b}\| \|\bm{b}\|} \\
&= \frac{k}{|k|} \\
&= \begin{cases}
+1 & \text{如果} \quad k > 0 \\
-1 & \text{如果} \quad k < 0
\end{cases}
\end{align}

因此$\theta = 0$或$\theta = \pi$，即$\bm{a}$和$\bm{b}$平行。

\section{应用到角动量和角速度}

\subsection{情况：绕主惯性轴旋转}

当角速度$\bm{\omega}$沿着主惯性轴方向时，我们有：
\begin{equation}
\bm{L} = \lambda_i \bm{\omega}
\end{equation}

其中$\lambda_i$是主转动惯量（标量）。

\subsection{应用向量平行的定义}

根据向量平行的定义，由于存在标量$\lambda_i$使得$\bm{L} = \lambda_i \bm{\omega}$，因此$\bm{L}$和$\bm{\omega}$平行。

\subsection{物理意义}

\begin{itemize}
\item $\lambda_i > 0$：角动量和角速度方向相同
\item $\lambda_i$的大小决定了角动量的大小相对于角速度的大小
\item 这是最"自然"的旋转方式，不需要额外的力矩来维持旋转
\end{itemize}

\section{反例：非平行的情况}

\subsection{一般情况}

在一般情况下，角动量$\bm{L} = I_b \bm{\omega}$，其中$I_b$是转动惯量矩阵。

\textbf{如果$I_b$不是标量矩阵}，那么$\bm{L}$和$\bm{\omega}$通常不平行。

\textbf{例子}：考虑转动惯量矩阵：
\begin{equation}
I_b = \begin{bmatrix}
I_{xx} & I_{xy} & 0 \\
I_{xy} & I_{yy} & 0 \\
0 & 0 & I_{zz}
\end{bmatrix}
\end{equation}

角速度$\bm{\omega} = \begin{bmatrix} \omega_x \\ \omega_y \\ 0 \end{bmatrix}$，则：
\begin{equation}
\bm{L} = \begin{bmatrix}
I_{xx} \omega_x + I_{xy} \omega_y \\
I_{xy} \omega_x + I_{yy} \omega_y \\
0
\end{bmatrix}
\end{equation}

除非$I_{xy} = 0$或$\omega_x = 0$或$\omega_y = 0$，否则$\bm{L}$和$\bm{\omega}$不平行。

\subsection{为什么不是标量倍数？}

如果$\bm{L} = k \bm{\omega}$，那么：
\begin{equation}
\begin{bmatrix}
I_{xx} \omega_x + I_{xy} \omega_y \\
I_{xy} \omega_x + I_{yy} \omega_y \\
0
\end{bmatrix} = k \begin{bmatrix}
\omega_x \\ \omega_y \\ 0
\end{bmatrix}
\end{equation}

这要求：
\begin{align}
I_{xx} \omega_x + I_{xy} \omega_y &= k \omega_x \\
I_{xy} \omega_x + I_{yy} \omega_y &= k \omega_y
\end{align}

对于一般的$\omega_x$和$\omega_y$，这只有在$I_{xy} = 0$且$I_{xx} = I_{yy} = k$时才能成立。

\section{向量平行的其他性质}

\subsection{性质1：传递性}

如果$\bm{a} \parallel \bm{b}$且$\bm{b} \parallel \bm{c}$，那么$\bm{a} \parallel \bm{c}$。

\textbf{证明}：如果$\bm{a} = k_1 \bm{b}$且$\bm{b} = k_2 \bm{c}$，那么$\bm{a} = k_1 k_2 \bm{c}$。

\subsection{性质2：标量乘法的保持}

如果$\bm{a} \parallel \bm{b}$，那么对于任意标量$k$，$k \bm{a} \parallel \bm{b}$。

\textbf{证明}：如果$\bm{a} = k_1 \bm{b}$，那么$k \bm{a} = k k_1 \bm{b}$。

\subsection{性质3：线性组合}

如果$\bm{a} \parallel \bm{b}$，那么$\bm{a} + \bm{b} \parallel \bm{b}$。

\textbf{证明}：如果$\bm{a} = k \bm{b}$，那么$\bm{a} + \bm{b} = (k + 1) \bm{b}$。

\section{总结}

\subsection{主要结论}

\begin{enumerate}
\item \textbf{定义}：两个向量平行当且仅当存在标量$k$使得$\bm{a} = k \bm{b}$。

\item \textbf{应用到角动量}：如果$\bm{L} = \lambda_i \bm{\omega}$（其中$\lambda_i$是标量），那么$\bm{L}$和$\bm{\omega}$平行。

\item \textbf{等价条件}：
\begin{itemize}
\item $\bm{a} = k \bm{b}$（标量倍数）
\item $\bm{a} \times \bm{b} = \bm{0}$（叉积为零）
\item $\cos\theta = \pm 1$（夹角为$0$或$\pi$）
\end{itemize}
\end{enumerate}

\subsection{关键要点}

\begin{itemize}
\item 标量倍数关系是向量平行的定义
\item 标量只改变向量的长度，不改变方向
\item 如果两个向量是标量倍数关系，它们必然平行
\item 这是数学上的定义，不是需要证明的定理
\end{itemize}

\subsection{应用场景}

\begin{itemize}
\item 刚体动力学中角动量与角速度的关系
\item 主惯性轴旋转的分析
\item 向量几何和线性代数中的基本概念
\end{itemize}

\end{document}

